\loai{Xác định I từ phương trình}
\begin{vidu}%[3P3-2.2B][Loai4] 
	\nguon{QG - 2020} Cường độ dòng điện $i=5\sqrt{2}\cos (100\pi t+\pi)\ (A)$ có giá trị hiệu dụng là
	\choice
	{$100\pi\ A$}
	{\True $5\ A$}
	{$\pi\ A$}
	{$5\sqrt{2}\ A$}	
	\loigiai{
		\begin{itemize}
			\item 
		\end{itemize}	
	}
\end{vidu}
\loai{Xác định T, f, $\varphi$ và pha từ phương trình}
\begin{vidu}%[3P3-2.2B][Loai4]
	\nguon{QG - 2020} Dòng điện xoay chiều chạy qua một đoạn mạch có cường độ $i=3\cos(100\pi t+\dfrac{\pi}{2})$ (A). Pha ban đầu của cường độ dòng điện này là
	\choice
	{3 rad}
	{\True $\dfrac{\pi}{2}$ rad}
	{$100\pi$ rad}
	{100 rad}
	\loigiai{Pha ban đầu của dòng điện là $\pi/2$}	
\end{vidu}
\begin{vidu} %[3P3-2.2B][Loai4] 
	Dòng điện xoay chiều có cường độ $i=2\sin50\pi t\ A$. Dòng điện này có
	\choice
	{cường độ hiệu dụng là $2\sqrt{2}\ A$}
	{\True tần số là $25\ Hz$}
	{cường độ hiệu dụng là $2\ A$}
	{chu kì là $0,02\ s$}
	\loigiai{
		\begin{itemize}
			\item Tần số: $f=\dfrac{\omega }{2\pi }=$$\dfrac{50\pi }{2\pi }=25\ Hz$.
		\end{itemize}
	}
\end{vidu}
\loai{Xác định giá trị tức thời tại thời điểm t}
\begin{vidu}%[3P3-2.2B][Loai6]
	Một điện áp xoay chiều có biểu thức $u=220\sqrt{2}\cos  \left( 100\pi t+\dfrac{2\pi}{3} \right)$ V (t tính bằng s). Điện áp tức thời tại $t = 0$ có giá trị
	\choice
	{$-110\sqrt{2}$ V và đang tăng}
	{\True $-110\sqrt{2}$ V  và đang giảm}
	{$110\sqrt{2}$ V  và đang giảm}
	{$110\sqrt{2}$ V và đang tăng}
	\loigiai{
		Tại $t = 0$: $\varphi =\dfrac{2\pi}{3}$
		$\Rightarrow$ $u=-\dfrac{U_0}{2}=-110\sqrt{2}$  và đang giảm.
	}
\end{vidu}
\loai{Xác định khoảng thời gian, thời điểm theo cường độ dòng điện}
\begin{vidu}%[3P3-2.2K][Loai7]
	Dòng điện xoay chiều qua một đoạn mạch có biểu thức $i={I_0}\cos \left( 120\pi t-\dfrac{\pi}{3} \right)\ A$. Thời điểm thứ 2018 độ lớn cường độ dòng điện bằng cường độ dòng điện hiệu dụng là
	\choice
	{8,15 s}
	{\True 8,4 s}
	{9,26 s}
	{10,3 s}
	\loigiai{
%\Binhluan{
		\begin{itemize}
			\item Tại t = 0: $i=\dfrac{I_0}{2}$ và đang tăng.
			\item $\left| i \right|=I=\dfrac{I_0}{\sqrt{2}}\Rightarrow i=\pm \dfrac{I_0}{\sqrt{2}}$, mỗi chu kì có 4 lần mà dòng điện $i=\pm \dfrac{I_0}{\sqrt{2}}$.
		\end{itemize}
%	}
%	{Có thể dùng VTLG cho dễ hình dung}
		\begin{itemize}	
			\centerline{\begin{tikzpicture}[scale=1,thick,>=stealth']
				\tkzDefPoints{0/0/o, 5/0/i, -5/0/i', 4/0/m, -4/0/n, 2/0/p, 2.8/0/q}
				\tkzDrawSegments[->](i',i)
				\tkzDrawPoints[size=3,fill=white](o,m,n,p,q)
				\tkzLabelPoint[above](o){$O$}
				\tkzLabelPoint[above](m){$I_0$}
				\tkzLabelPoint[above](p){$\dfrac{I_0}{2}$}
				\tkzLabelPoint[above](n){$-I_0$}
				\tkzLabelPoint[above](q){$\dfrac{I_0\sqrt{2}}{2}$}
				\tkzLabelPoint[above](i){$i$}
				\draw[->](2,-0.5)--(4,-0.5)--++(-90:0.3)--(2.8,-0.8);
				\end{tikzpicture}}			
			\item Tách $2018 = 2016 + 2$.
			Thời điểm cần tìm là: ${{t}}=504T+\dfrac{T}{6}+\dfrac{T}{8}\approx 8,4\ s$.
	\end{itemize}
}
\end{vidu}
\loai{Xác định khoảng thời gian, thời điểm theo điện áp}
\begin{vidu}%[3P3-2.2K][Loai8]
	Hiệu điện thế giữa hai đầu một đoạn mạch có biểu thức $u = U_0\sin(100\pi t + \dfrac{\pi}{2})$  V. Tại thời điểm t nào sau đây hiệu điện thế tức thời $u \ne \dfrac{U_0}{\sqrt 2}$?
	\choice
	{$\dfrac{1}{400}$ s}
	{$\dfrac{9}{400}$ s}
	{$\dfrac{7}{400}$ s}
	{\True $\dfrac{11}{400}$ s}
	\loigiai{
%		\Binhluan{
		Ta có: 
		\begin{align*}
		&u = U_0\sin(100\pi t + \dfrac{\pi}{2})= {U_0}\cos  100\pi t = \dfrac{U_0}{\sqrt 2}\\ 
		\Rightarrow\ &\begin{cases} 100\pi t = \dfrac{\pi}{4} + 2k\pi \Rightarrow t = \dfrac{1 + 8k}{400}\ s\\
		100\pi t = - \dfrac{\pi}{4} + 2k\pi \Rightarrow t = \dfrac{8k - 1}{400}\ s \end{cases}.
	\end{align*}
%		}
%	{Chuẩn hoá $U_0=100\ V$ rồi bấm máy tính cũng là 1 ý hay}
	}
\end{vidu}
\loai{Xác định u, i tại thời điểm sau khi biết tại thời điểm trước}
\begin{vidu} %[3P3-2.2K][Loai10] 
	\nguon{CĐ - 2013} Điện áp ở hai đầu một đoạn mạch là $u=160cos100\pi t$ V (t tính bằng giây). Tại thời điểm $t_1$, điện áp ở hai đầu đoạn mạch có giá trị là 80\ V và đang giảm. đến thời điểm $t_2=t_1+0,015$ s, điện áp ở hai đầu đoạn mạch có giá trị bằng
	\choice
	{$40\sqrt{3}$ V}
	{\True $80\sqrt{3}$ V}
	{40 V}
	{80 V}
	\loigiai{
		immini[thm]{\begin{itemize}
				\item $t_2=t_1+0,015\ s = t_1+ \dfrac{3T}{4}$.
				\item Với $\dfrac{3T}{4}$ ứng góc quay $\dfrac{3\pi}{2}$. 
				\item Từ $t_1$ đến $t_2$ góc quay $M_1OM_2 =\dfrac{3\pi}{2}$.
				\item Nhìn hình vẽ thời gian quay $\dfrac{3T}{4}$ (ứng với góc quay $\dfrac{3\pi}{2}$).
				\item $M_2$ chiếu xuống trục u $\Rightarrow u=80\sqrt{3}$ V. 
		\end{itemize}}{\begin{tikzpicture}[scale=1,thick,>=stealth']
			\tkzInit[ymin=-2.5,ymax=2.5,xmin=-3,xmax=3]\tkzClip
			\tkzDefPoints{-2.5/0/m, 2.5/0/n, 0/0/o, 0/2.5/p, 0/-2.5/q}
			\tkzDefShiftPoint[o](-30:2){m1}
			\tkzDefShiftPoint[o](60:2){m2}
			\tkzDefPointBy[projection=onto m--n](m1) \tkzGetPoint{h1}
			\tkzDefPointBy[projection=onto m--n](m2) \tkzGetPoint{h2}
			\draw(o)circle(2cm);
			\tkzDrawSegments[dashed,thin](m1,h1 m2,h2)
			\tkzDrawSegments[->](m,n)
			\tkzDrawSegments(p,q)
			\tkzDrawSegments[blue,->](o,m1 o,m2)
			\tkzDrawPoints[size=3,fill=white](o,m1,m2,h1,h2)
			\tkzMarkAngles[size=0.4cm](n,o,m2)
			\tkzLabelAngles[pos=0.7](n,o,m2){$\dfrac{\pi}{3}$}
			\node at (m1)[below right]{$M_2$};
			\node at (m2)[above right]{$M_1$};
			\node at (2,0)[below right]{$160$};
			\node at (n)[above]{$u$};
			\node at (h1)[above]{$80\sqrt{3}$};
			\node at (h2)[below]{$80$};
			\end{tikzpicture}}
	}
\end{vidu}
\loai{Xác định thời gian đèn sáng hoặc tắt}
\begin{vidu}%[3P3-2.2B][Loai11]
	Một đèn ống được mắc vào mạng điện xoay chiều có phương trình $u = 220\sqrt{2}\cos \left( 100\pi t-\dfrac{\pi}{2} \right)$ V (trong đó u tính bằng V, t tính bằng s). Biết rằng đèn sáng mỗi khi điện áp hai đầu đèn có độ lớn không nhỏ hơn $110\sqrt{2}\ V$. Khoảng thời gian đèn tắt trong một chu kì là
	\choice
	{$\dfrac{1}{300}\ s$}
	{\True $\dfrac{1}{150}\ s$}
	{$\dfrac{1}{75}\ s$}
	{$\dfrac{1}{50}\ s$}
	\loigiai{
		Mỗi chu kì, khoảng thời gian đèn tắt (nét đứt) là $\Delta t=\dfrac{T}{3}=\dfrac{1}{150}\ s$.
		\\
		\centerline{\begin{tikzpicture}[scale=1,thick,>=stealth']
			\tkzDefPoints{0/0/o, 5/0/i, -5/0/i', 4/0/m, -4/0/n, 2/0/p, -2/0/p', 3.5/0/q}
			\tkzDrawSegments[->](i',i)
			\tkzDrawPoints[size=3,fill=white](o,m,n,p,p')
			\tkzLabelPoint[above](o){$O$}
			\tkzLabelPoint[above](m){$U_0$}
			\tkzLabelPoint[above](p){$\dfrac{U_0}{2}$}
			\tkzLabelPoint[above](p'){$-\dfrac{U_0}{2}$}
			\tkzLabelPoint[above](n){$-U_0$}
			\tkzLabelPoint[above](i){$u$}
			\draw[very thick](2,-0.5)--(4,-0.5)--++(-90:0.3)--(2,-0.8);
			\draw[very thick](-2,-0.5)--(-4,-0.5)--++(-90:0.3)--(-2,-0.8);
			\draw[dashed](-2,-0.5)--(2,-0.5)(-2,-0.8)--(2,-0.8);
			\end{tikzpicture}}		
	}
\end{vidu}
\loai{Xác định thời điểm đèn sáng hoặc tắt}
\begin{vidu}%[3P3-2.2K][Loai12]
	Đặt vào hai đầu đèn ống điện áp xoay chiều $u=220\sqrt{2}\cos \left( \dfrac{100\pi t}{3}+\dfrac{\pi}{2} \right)$ V. Biết đèn chỉ sáng khi điện áp tức thời có độ lớn không nhỏ hơn $110\sqrt{2}$ V. Kể từ $t = 0$, thời điểm đèn sáng lần thứ 2018 là
	\choice
	{60,505 s}
	{\True 60,515 s}
	{30,275 s}
	{30,265 s}
	\loigiai{
		\begin{itemize}
			\item Tại $t = 0:\ u = 0$  và đang giảm.
			\\
			\centerline{\begin{tikzpicture}[scale=1,thick,>=stealth']
				\tkzDefPoints{0/0/o, 5/0/i, -5/0/i', 4/0/m, -4/0/n, 2/0/p, -2/0/p', 3.5/0/q}
				\tkzDrawSegments[->](i',i)
				\tkzDrawPoints[size=3,fill=white](o,m,n,p,p')
				\tkzLabelPoint[above](o){$O$}
				\tkzLabelPoint[above](m){$U_0$}
				\tkzLabelPoint[above](p){$\dfrac{U_0}{2}$}
				\tkzLabelPoint[above](p'){$-\dfrac{U_0}{2}$}
				\tkzLabelPoint[above](n){$-U_0$}
				\tkzLabelPoint[above](i){$u$}
				\draw[->](0,-0.5)--(-4,-0.5)--++(-90:0.3)--(2,-0.8);
				\end{tikzpicture}}			
			\item Mỗi chu kì, đèn sáng 2 lần $\Rightarrow$ tách $2018 = 2016 + 2 \Rightarrow$ $\Delta t = 1008T +$ $\dfrac{T}{4}+\dfrac{T}{3} = 60,515$ s.
	\end{itemize}}
\end{vidu}
\loai{Xác định giá trị cực đại và hiệu dụng từ đồ thị}
\begin{vidu}%[3P3-2.2B][Loai13]
	immini[thm]{Đồ thị phụ thuộc thời gian của cường độ dòng điện chạy qua mạch như hình vẽ. Cường độ hiệu dụng là
		\motcot
		{3A}
		{3,5 A}
		{5 A}
		{\True 2,5$\sqrt{2}$ A}}{\begin{tikzpicture}[scale=0.75,thick,>=stealth',x=1.5cm,y=1cm]
		\draw[->] (0,0)--(3.5,0)node[above]{$t$};
		\draw[->] (0,-2)--(0,2)node[right]{$i\ (A)$};
		
		\foreach  \A  in  {1.5} % Biên độ
		\foreach  \T  in  {3} % Chu kì
		\foreach  \Phi  in  {-90} % Pha ban đầu
		{\draw [red, line width = 1.2pt, domain=0:3.4, samples=500]%
			plot (\x, {(\A)*cos(360/\T*\x+\Phi)});
			\draw[dashed](0, {\A})node[left]{$5$}--(0.75, {(\A)*cos(360/\T*0.75+\Phi)})--(0.75,0);
		}
		\node at (0,0)[left]{$O$};
		\draw[fill=white](0.75,1.5)circle(1.2pt);
		\end{tikzpicture}}
	\loigiai{
		Theo đồ thị ta có: $I_0 = 5\ A \Rightarrow I = \dfrac{I_0}{\sqrt{2}} = 2,5\sqrt{2}$ A.
	}
\end{vidu}
\loai{Viết biểu thức u, i từ đồ thị}
\begin{vidu}%[3P3-2.2B][Loai14]
	immini[thm]{Trên hình vẽ là đồ thị phụ thuộc thời gian của dòng điện xoay chiều. Biểu thức của dòng điện là
		\motcot
		{$i = 2\cos(50\pi t + \pi)$ A}
		{\True $i = 2\cos(50\pi t - \dfrac{\pi}{3})$ A}
		{$i = 2\cos(50\pi t + \dfrac{\pi}{3})$ A}
		{$i = 2\cos(50\pi t)$ A}}{\begin{tikzpicture}[scale=1,thick,>=stealth',x=1cm,y=1cm]
		\draw[->] (0,0)--(3,0)node[above,xshift=7pt]{$t\ (s)$};
		\draw[->] (0,-2)--(0,2)node[right]{$i\ (A)$};
		
		\foreach  \A  in  {1.5} % Biên độ
		\foreach  \T  in  {2.4} % Chu kì
		\foreach  \Phi  in  {-60} % Pha ban đầu
		{\draw [red, line width = 1.2pt, domain=0:2.5, samples=500]%
			plot (\x, {(\A)*cos(360/\T*\x+\Phi)});
		}
		\node at (0,0)[left]{$O$};
		\draw[dashed](0,1.5)--(0.4,1.5);
		\draw[fill=white](0,1.5)circle(1.2pt);
		\draw[fill=white](0,0.75)circle(1.2pt);
		\draw[fill=white](1,0)circle(1.2pt);
		\draw[fill=white](0.4,1.5)circle(1.2pt);
		\node at (0,1.5)[left]{$2$};
		\node at (0,0.75)[left]{$1$};
		\node at (1,0)[below left]{$\dfrac{1}{60}$};
		\end{tikzpicture}}
	\loigiai{
		Tại t = 0 thì i = $\dfrac{I_0}{2}$ và đang tăng $\Rightarrow \varphi  = - \dfrac{\pi}{3}$.
	}
\end{vidu}
\loai{Xác định độ lệch pha từ đồ thị}
\begin{vidu}%[3P3-2.2K][Loai15]
	immini[thm]{Hình vẽ bên là đồ thị phụ thuộc thời gian của hai dòng điện xoay chiều $(1)$ và $(2)$. So với dòng điện $(1)$ thì dòng điện $(2)$
		\motcot
		{sớm pha hơn $\dfrac{\pi}{12}$}
		{\True sớm pha hơn $\dfrac{\pi}{6}$}
		{trễ pha hơn $\dfrac{\pi}{6}$}
		{trễ pha hơn $\dfrac{\pi}{12}$}}{\begin{tikzpicture}[scale=1,thick,>=stealth',x=1.5cm,y=1cm]
		\draw[->] (0,0)--(3.2,0)node[above]{$t\ (s)$};
		\draw[->] (0,-2)--(0,2)node[right]{$i\ (A)$};
		\foreach  \y  in  {-1.5,1.5}
		{\draw [dashed,thin] (0,\y) --(3,\y);}
		\foreach  \x  in  {0.2,0.4,...,3}
		{\draw [dashed,thin] (\x,-1.5) --(\x,1.5);}
		
		\foreach  \A  in  {1.5} % Biên độ
		\foreach  \T  in  {2.4} % Chu kì
		\foreach  \Phi  in  {60} % Pha ban đầu
		{\draw [red, line width = 1.2pt, domain=0:3, samples=500]%
			plot (\x, {(\A)*cos(360/\T*\x+\Phi)});
		}
		
		\foreach  \A  in  {0.75} % Biên độ
		\foreach  \T  in  {2.4} % Chu kì
		\foreach  \Phi  in  {30} % Pha ban đầu
		{\draw [blue,dashed, line width = 1.2pt, domain=0:3, samples=500]%
			plot (\x, {(\A)*cos(360/\T*\x+\Phi)});
		}
		
		\node at (0,0)[left]{$O$};
		\node at (3,-1)[below]{$(2)$};
		\node at (3,-0.2)[below]{$(1)$};
		\end{tikzpicture}}
	\loigiai{
	\Binhluan{
		\begin{itemize}
			\item Chu kì $T = 12$ ô $\sim 2\pi$.
			\item $i_2$ cắt trục hoành sớm hơn $i_1$ 1 ô $\sim \dfrac{\pi}{6}$.
	\end{itemize}
}
{Dùng VTLG dễ thấy hơn}
}
\end{vidu}
